\section{Introduction}
\label{sec:intro}

Language models are increasingly asked to describe sensor data: activity from wearables, driving from phones, machines from their vibration~\cite{xu2024penetrative,li2025sensorllm,zhang2025sensorlm}.
The data-to-text answer to the question of trust is to compile the measurements into facts, generate text from them, and verify every claim against the facts, falling back to a deterministic narrator when verification fails~\cite{wiseman2017challenges,moryossef2019step,jhamtani2021truce}.
Such a system is \emph{faithful}: it never says what its facts do not say~\cite{maynez2020faithfulness}.

This paper is about what faithfulness leaves open.
A verifier compares a sentence with a value; it cannot tell whether the value determines what the sentence asserts.
We call a claim \emph{identifiable} when it does, and study two ways in which it fails.
At the \emph{world level}, the value is not a property of what the sentence talks about. The vibration a car records on a road depends on the car and its speed, so ``the surface is in fair condition'' can agree with one car's ledger and contradict another car's ledger for the same street.
At the \emph{input level}, a fine-tuned narrator is trained on targets that state values its input does not contain. It learns to produce such values without evidence, and the verifier then rejects its output: safely, but the model contributes nothing.
We study both in RoadScribe, a verified narrator for IMU and GNSS recordings of car trips. We use the PVS dataset, in which three vehicles drove the same three routes, and a 3B-parameter language model fine-tuned under controlled conditions.

\paragraph{Contributions.}
\begin{enumerate}
\item \emph{Two-level claim identifiability, with an operational test for each level} (Sect.~\ref{sec:ident}). The world-level test compares cross-condition agreement with a same-condition noise floor and a placebo. The input-level test is a typed audit of whether each value the training targets state is present in the input. Both are cheap, model-independent and applicable to other sensor-to-text systems.
\item \emph{World-level evidence, pre-registered} (Sect.~\ref{sec:world}).
  \begin{itemize}
  \item RoadScribe's road class differs between two vehicles at the same place in 76\% of comparisons, and the vehicle explains 18\% of the variance of the road value. Two sensors in one vehicle agree at $\kappa=0.98$.
  \item A speed elasticity estimated across places has the wrong sign.
  \item One gain per vehicle, estimated on about a kilometre of shared road, raises cross-vehicle agreement from $\kappa=0.72$ to $0.96$, and the accuracy on road-quality labels of held-out routes and vehicles from 90.5\% to 98.3\%. The result replicates on a second logger.
  \end{itemize}
\item \emph{Input-level evidence from a controlled experiment, pre-registered} (Sect.~\ref{sec:input}). \PSIXCONTRIB
\item \emph{A design consequence.} When the model only selects content and code writes the values, every valid model plan yields a verified narrative.
\end{enumerate}

\paragraph{Scope.}
We propose no new roughness estimator, language model or verifier. Per-vehicle calibration is established practice~\cite{sayers1986iri}.
Our claims concern which statements a sensor-to-text system may make, and how to test this before deployment.
They rest on one public dataset (three vehicle--driver pairs, one pass per route, fixed modules rather than phones) and on a synthetic training corpus.
